% !TeX root = ../../Book3.tex
\chapter{\texorpdfstring{\(\theta\)}{\AnalysisPdfTheta} 函数与整数平方和}
\label{b3:ch:24}
\BookChapterNumber{b3:ch:24}
% 本章摘要（不计入编号）
\begin{chapterabstract}
theta 级数把高斯求和、复函数的零点和整数表示数放在同一个函数中。本章先从收敛与指标替换建立四种 theta 函数、热方程及模变换，再以有限恒等式的极限证明 Jacobi 三重积。乘积确定零点和导数常数，椭圆函数的半周期值将这些资料转成四平方和的约数公式。二平方和另由 Gaussian 整数的除法与唯一分解得到。末节说明一般格的反演律包含对偶格和体积，并在偶自对偶条件下回到前章的全模群模形式。
\end{chapterabstract}

% 本章目标（不计入编号）
\begin{learninggoals}
能够核对 \(q=e^{\pi\mathrm{i}\tau}\) 与前章变量的区别，说明级数求导、乘积取极限和系数重排的依据；能够从乘积找出 theta 零点，从表示数的定义重建生成函数；能够说明二平方和的素因子障碍，并把半周期差、热方程、Lambert 级数连接成四平方和证明。选读后还应能辨别一般格的反演公式与全模群模形式的条件，并计算八维偶自对偶格的短向量数。
\end{learninggoals}

把
\[
1+2q+2q^4+2q^9+\cdots
=\sum_{n\in\mathbb Z}q^{n^2}
\]
平方，\(q^m\) 的系数就在数两个平方相加等于 \(m\) 的方式。这里系数二记录正负号；例如 \(q^1\) 来自 \(1^2\) 与 \((-1)^2\)，而常数项只有零的一种选择。这个生成函数容易写出，但把它相乘并不能立刻给出所有系数的公式。需要找到一种表达式，让整数 \(m\) 的约数或素因子能够显露出来。

本章在三个表示之间来回转换：级数便于计数和微分，乘积便于判断零点和取对数导数，变换公式则把一个参数区域与另一个区域联系起来。每次转换都有具体用途。Poisson 求和解释高斯参数的反演；三重积使 theta 常数的非零性变得可见；椭圆函数的主部比较再把一个对数导数变成 theta 常数的四次幂。最后一步展开几何级数，才得到约数求和。

\ProseParagraphBreak
阅读本章时，只需从前面带来少量明确工具：正规收敛允许逐项求导，正规乘积的非零因子给出非零极限；\chapref{b3:ch:22}中无极点的椭圆函数为常数，Weierstrass 函数在原点有展开 \(\wp(z)=z^{-2}+O(z^2)\)。后面真正使用这些结论时会再次指出其作用。所需的高斯求和在本章重建；二平方和需要的 Gaussian 整数分解也完整证明，不预设代数数论背景。末节的偶自对偶格例子才进一步调用\chapref{b3:ch:23}的模形式空间。

\secref{b3:sec:2401}至\secref{b3:sec:2404}构成基本路线；\secref{b3:sec:2405}将同一方法推广到高维格，可作为选读。Whittaker、Watson 的《A Course of Modern Analysis》\cite{Whittaker1996Course}可供查阅 theta 与椭圆函数的不同记号；Zagier 的《Elliptic Modular Forms and Their Applications》\cite[pp. 26--28]{Zagier2008Elliptic}还说明用模形式空间比较平方和系数的方法。阅读外部公式时，应先核对空间变量的周期以及 \(q\) 的定义，再比较数值常数。

% !TeX root = ../../Book3.tex
\section{Jacobi \texorpdfstring{\(\theta\)}{\AnalysisPdfTheta} 函数}
\label{b3:sec:2401}

若要用一个函数同时保存所有平方数，最直接的办法是把平方数放到幂指数上，写成 \(\sum_{n\in\mathbb Z}q^{n^2}\)。当 \(|q|<1\) 时，这个级数迅速收敛；乘起若干份后，系数就在计数平方和。它还有一个计数问题本身不容易看出的对称：令 \(q=e^{-\pi t}\)，把 \(t\) 换成 \(1/t\)，级数只改变一个简单因子。这一对称来自高斯函数的 Fourier 变换。下面加入一个记录线性相位的空间变量 \(z\)，把这项事实建立为可以求导、平移和作参数变换的函数恒等式。

\subsection{\texorpdfstring{\(\theta\)}{\AnalysisPdfTheta} 级数}
\label{b3:subsec:240101}

\begin{definition}\label{b3:def:24-theta-series}
设 \(\tau\in\mathfrak H=\{\,w\in\mathbb C\mid\operatorname{Im}w>0\,\}\)、\(z\in\mathbb C\)，置 \(q=e^{\pi\mathrm i\tau}\)。四个\term[jacobi-theta-hanshu]{Jacobi theta 函数}[Jacobi theta function]定义为
\begin{equation}\label{b3:eq:24-theta-series}
\begin{aligned}
 \theta_1(z,\tau)&=-\mathrm i\sum_{n\in\mathbb Z}(-1)^n
                q^{(n+1/2)^2}e^{(2n+1)\mathrm iz},\\
 \theta_2(z,\tau)&=\sum_{n\in\mathbb Z}
                q^{(n+1/2)^2}e^{(2n+1)\mathrm iz},\\
 \theta_3(z,\tau)&=\sum_{n\in\mathbb Z}q^{n^2}e^{2n\mathrm iz},\\
 \theta_4(z,\tau)&=\sum_{n\in\mathbb Z}(-1)^nq^{n^2}e^{2n\mathrm iz}.
\end{aligned}
\end{equation}
这里对任意实数 \(a\)，记号 \(q^a\) 一律表示 \(e^{\pi\mathrm i a\tau}\)，不另选复数 \(q\) 的幂分支。省略第二个变量时，始终保持同一个 \(\tau\)。
\end{definition}
\symidx{\(\theta_j(z,\tau)\)：Jacobi theta 函数，\(1\le j\le4\)}

本章的 \(q\) 与前章的模形式展开变量不同：前章用 \(e^{2\pi\mathrm i\tau}\)，本章用它的平方根 \(e^{\pi\mathrm i\tau}\)。这一选择使 \(\theta_3(0,\tau)\) 的幂指数直接成为整数平方，而不再带二分之一。对 \(\theta_1,\theta_2\)，指数还含四分之一；它们首先是 \(\tau\) 的单值函数，不能未经说明就把它们看成穿孔圆盘中 \(q\) 的单值函数。
\ProseParagraphBreak
四个函数不是四种互不相关的构造。由整数格改为半整数格，就从 \(\theta_3\) 得到 \(\theta_2\)；再给相邻格点加上交替符号，就出现 \(\theta_4\) 和 \(\theta_1\)。这些变化在 Fourier 变换下会彼此交换，所以把它们一并引入，可以让后面的变换律保持封闭。先把 \(n\) 与 \(-n\)、或 \(n\) 与 \(-n-1\) 配对，得到可直接计算的实三角形式：
\[
\begin{aligned}
 \theta_1(z,\tau)&=2\sum_{n\ge0}(-1)^nq^{(n+1/2)^2}
                                      \sin((2n+1)z),\\
 \theta_2(z,\tau)&=2\sum_{n\ge0}q^{(n+1/2)^2}
                                      \cos((2n+1)z),\\
 \theta_3(z,\tau)&=1+2\sum_{n\ge1}q^{n^2}\cos(2nz),\\
 \theta_4(z,\tau)&=1+2\sum_{n\ge1}(-1)^nq^{n^2}\cos(2nz).
\end{aligned}
\]
因此 \(\theta_1\) 关于 \(z\) 为奇函数，其余三个为偶函数；\(\theta_1(0,\tau)=0\)。这里的配对和随后使用的逐项求导，都由下面的同一估计保证。

\begin{proposition}\label{b3:prop:24-theta-normal}
对 \(j\in\{1,2,3,4\}\)，\cref{b3:def:24-theta-series}的级数及其任意固定阶混合导数
\(\partial_z^r\partial_\tau^s\)，其中 \(r,s\) 为非负整数，都在
\(\mathbb C\times\mathfrak H\) 的每个紧子集上绝对一致收敛。因而 \(\theta_j(z,\tau)\) 对 \(z\) 为整函数，对 \(\tau\in\mathfrak H\) 全纯，所述导数可以逐项计算。
\end{proposition}
\begin{proof}
在一个给定紧子集上，可以选取 \(\eta>0\)、\(R>0\)，使
\(\operatorname{Im}\tau\ge\eta\)、\(|\operatorname{Im}z|\le R\)。四种级数的指标都是 \(u=n\) 或 \(u=n+1/2\)，除去模为一的系数后，各项及导数的模不超过
\[
 C_{r,s}|u|^{r+2s}\exp(-\pi\eta u^2+2R|u|).
\]
当 \(|u|\) 足够大时，\(2R|u|\le\pi\eta u^2/2\)，故尾项被
\(C_{r,s}|u|^{r+2s}e^{-\pi\eta u^2/2}\) 控制。这个数列可求和，因为任意固定次多项式都被高斯衰减吸收。有限个小指标另取上界，Weierstrass 判别法便给出一致收敛；对单个变量应用全纯函数列的逐项求导定理，即得到各阶混合导数。
\end{proof}

上半平面条件确实参与了论证。若 \(\tau\) 为实数、\(z=0\)，则 \(\theta_3\) 的各项模恒为一，连级数收敛所需的项趋于零都不满足。不能把上面的紧集估计直接延到 \(\operatorname{Im}\tau=0\)。
\ProseParagraphBreak
平方指数还带来一个将在四平方和证明中使用的微分关系：对参数求导产生平方，对空间变量求两次导数也产生同一个平方。

\begin{proposition}\label{b3:prop:24-theta-heat}
对 \(z\in\mathbb C\)、\(\tau\in\mathfrak H\) 和 \(1\le j\le4\)，有
\begin{equation}\label{b3:eq:24-theta-heat}
 \frac1{\pi\mathrm i}\partial_\tau\theta_j(z,\tau)
       =-\frac14\partial_z^2\theta_j(z,\tau).
\end{equation}
若以 \(q=e^{\pi\mathrm i\tau}\) 作局部坐标，则左端为
\(q\partial_q\theta_j\)。特别地，函数 \(u_j(x,t)=\theta_j(x,\mathrm it)\)，
\(x\in\mathbb R\)、\(t>0\)，满足
\(\partial_tu_j=(\pi/4)\partial_x^2u_j\)。
\end{proposition}
\begin{proof}
每项都具有常数乘 \(e^{\pi\mathrm i u^2\tau+2\mathrm iuz}\) 的形式，
其中 \(u=n\) 或 \(n+1/2\)。算子 \((\pi\mathrm i)^{-1}\partial_\tau\)
乘出 \(u^2\)，而 \(\partial_z^2\) 乘出 \(-4u^2\)。由
\cref{b3:prop:24-theta-normal}逐项计算并求和，得到恒等式。
由于 \(\mathrm dq/\mathrm d\tau=\pi\mathrm i q\ne0\)，局部坐标中的链式法则给出 \(q\partial_q=(\pi\mathrm i)^{-1}\partial_\tau\)。
最后 \(\partial_tu_j=\mathrm i\partial_\tau\theta_j\)，代入即得实时间的热方程。
\end{proof}

\begin{proposition}\label{b3:prop:24-theta-translations}
设 \(\tau\in\mathfrak H\)、\(z\in\mathbb C\)，并在下式中简记
\(\theta_j(z)=\theta_j(z,\tau)\)、\(q^a=e^{\pi\mathrm i a\tau}\)。
四个 theta 函数满足
\begin{equation}\label{b3:eq:24-theta-quasiperiods}
\begin{aligned}
 \theta_1(z+\pi)&=-\theta_1(z),&
 \theta_1(z+\pi\tau)&=-q^{-1}e^{-2\mathrm iz}\theta_1(z),\\
 \theta_2(z+\pi)&=-\theta_2(z),&
 \theta_2(z+\pi\tau)&=q^{-1}e^{-2\mathrm iz}\theta_2(z),\\
 \theta_3(z+\pi)&=\theta_3(z),&
 \theta_3(z+\pi\tau)&=q^{-1}e^{-2\mathrm iz}\theta_3(z),\\
 \theta_4(z+\pi)&=\theta_4(z),&
 \theta_4(z+\pi\tau)&=-q^{-1}e^{-2\mathrm iz}\theta_4(z).
\end{aligned}
\end{equation}
半周期平移的结果由\cref{b3:tab:24-half-translations}给出。
\end{proposition}

\begin{table}[htbp]
\centering
\caption[theta 函数的半周期平移]{theta 函数的半周期平移。共同参数为
\(\tau\in\mathfrak H\)，\(q=e^{\pi\mathrm i\tau}\)，每一行对全部 \(z\in\mathbb C\) 成立。第三列已乘去共同指数因子，因而只留下函数交换及一个常数。}
\label{b3:tab:24-half-translations}
\begin{tabular}{ccc}
\toprule
 \(j\) & \(\theta_j(z+\pi/2,\tau)\) &
 \(q^{1/4}e^{\mathrm iz}\theta_j(z+\pi\tau/2,\tau)\)\\
\midrule
 \(1\) & \(\theta_2(z,\tau)\) & \(\mathrm i\theta_4(z,\tau)\)\\
 \(2\) & \(-\theta_1(z,\tau)\) & \(\theta_3(z,\tau)\)\\
 \(3\) & \(\theta_4(z,\tau)\) & \(\theta_2(z,\tau)\)\\
 \(4\) & \(\theta_3(z,\tau)\) & \(\mathrm i\theta_1(z,\tau)\)\\
\bottomrule
\end{tabular}
\end{table}
\begin{proof}
对 \(z\mapsto z+\pi\)，整数频率的项不变，半整数频率的项乘以 \(-1\)，给出第一列公式。对 \(z\mapsto z+\pi\tau\)，把指数
\(u^2+2u\) 写成 \((u+1)^2-1\)，再把指标平移一位，提出共同因子 \(q^{-1}e^{-2\mathrm iz}\)。带 \((-1)^n\) 的级数在移指标后另产生一个负号，故得到第二列。
半周期表的第二列来自
\(e^{2\mathrm in(\pi/2)}=(-1)^n\) 与
\(e^{(2n+1)\mathrm i\pi/2}=\mathrm i(-1)^n\)。
对第三列，整数指标满足
\(n^2+n=(n+1/2)^2-1/4\)，半整数指标满足
\((n+1/2)^2+(n+1/2)=(n+1)^2-1/4\)。
例如对奇函数作这个移项，得到
\[
\begin{aligned}
 \theta_1(z+\pi\tau/2,\tau)
 &=-\mathrm i q^{-1/4}
       \sum_{n\in\mathbb Z}(-1)^nq^{(n+1)^2}e^{(2n+1)\mathrm iz}\\
 &=\mathrm i q^{-1/4}e^{-\mathrm iz}\theta_4(z,\tau).
\end{aligned}
\]
去掉交替符号就得到
\(\theta_2(z+\pi\tau/2,\tau)=q^{-1/4}e^{-\mathrm iz}\theta_3(z,\tau)\)。
整数指标的两行直接用第一项配方，分别得到
\(q^{-1/4}e^{-\mathrm iz}\theta_2(z,\tau)\) 与
\(\mathrm i q^{-1/4}e^{-\mathrm iz}\theta_1(z,\tau)\)，完成全表。
所有移项均在绝对收敛级数中进行。
\end{proof}

第二个方向的平移带有依赖 \(z\) 的乘法因子，这称为准周期性，不能把 \(\pi\tau\) 称作真正周期。与第二十二章的椭圆函数相比，这个差别很必要：非恒定整函数若有两个实线性无关的真正周期，就会在一个紧基本平行四边形上有界，继而在全平面有界，与 Liouville 定理矛盾。theta 函数允许大小在相邻格胞间按指数因子改变，因而能够同时保持整性与非平凡的格对称。取一次对数导数，乘法指数变成加法常数；再求导一次，常数消失，才会产生真正双周期的亚纯函数。

\subsection{与第二卷 Poisson 求和的联系}
\label{b3:subsec:240102}

第二卷从 Fourier 理论研究过 Poisson 求和。本章只需要高斯及其平移、调制的情形，下面直接证明这个版本，并把积分换序和 Fourier 展开成立的理由写出。这样，即使读者忘记了一般定理的函数空间条件，也可以沿着当前论证继续学习。

\begin{lemma}\label{b3:lem:24-gaussian-poisson}
设 \(t>0\)、\(a\in\mathbb R\)、\(b\in\mathbb C\)。则下列两边级数绝对收敛，并且
\begin{equation}\label{b3:eq:24-gaussian-poisson}
 \sum_{n\in\mathbb Z}e^{-\pi t(n+a)^2}e^{2\pi\mathrm ib(n+a)}
 =t^{-1/2}\sum_{m\in\mathbb Z}
       e^{-\pi(m-b)^2/t}e^{2\pi\mathrm ima}.
\end{equation}
特别地，\(a=b=0\) 时有
\[
 \sum_{n\in\mathbb Z}e^{-\pi tn^2}
     =t^{-1/2}\sum_{m\in\mathbb Z}e^{-\pi m^2/t}.
\]
\end{lemma}
\begin{proof}
固定 Fourier 变换约定
\(\widehat f(\xi)=\int_{\mathbb R}f(x)e^{-2\pi\mathrm ix\xi}\,\mathrm dx\)。
先计算 \(g_t(x)=e^{-\pi tx^2}\) 的变换。高斯积分为
\(\int_{\mathbb R}e^{-\pi tx^2}\,\mathrm dx=t^{-1/2}\)：在 \(t=1\)
时将积分平方化成平面上的非负积分，极坐标给出
\(2\pi\int_0^\infty e^{-\pi r^2}r\,\mathrm dr=1\)，一般情形再缩放。
令 \(F(\xi)=\widehat g_t(\xi)\)，对实 \(\xi\) 在积分下求导，并用
\(g_t'(x)=-2\pi txg_t(x)\) 分部积分，得到
\[
 F'(\xi)=\frac{\mathrm i}{t}
       \int_{\mathbb R}g_t'(x)e^{-2\pi\mathrm ix\xi}\,\mathrm dx
       =-\frac{2\pi\xi}{t}F(\xi).
\]
这里 \(g_t\) 与 \(xg_t\) 可积，且 \(g_t(x)\to0\) 保证边界项消失。
初值 \(F(0)=t^{-1/2}\) 因而给出
\[
 \widehat g_t(\xi)=t^{-1/2}e^{-\pi\xi^2/t}.
\]
对实数 \(b\)，令 \(f_b(x)=g_t(x)e^{2\pi\mathrm ibx}\)，于是
\(\widehat f_b(\xi)=t^{-1/2}e^{-\pi(\xi-b)^2/t}\)。

将 \(f_b\) 周期化为
\(P_b(x)=\sum_{n\in\mathbb Z}f_b(x+n)\)。当 \(x\in[0,1]\) 时，
\(|x+n|\ge |n|-1\)，而 \(f_b\) 的各阶导数都是多项式乘高斯再乘一个模为一的相位。故级数及其各阶导数一致收敛，\(P_b\) 为光滑的一周期函数。其第 \(m\) 个 Fourier 系数是
\[
\begin{aligned}
 \int_0^1P_b(x)e^{-2\pi\mathrm imx}\,\mathrm dx
 &=\sum_{n\in\mathbb Z}\int_0^1f_b(x+n)e^{-2\pi\mathrm imx}\,\mathrm dx\\
 &=\int_{\mathbb R}f_b(u)e^{-2\pi\mathrm imu}\,\mathrm du
   =t^{-1/2}e^{-\pi(m-b)^2/t}.
\end{aligned}
\]
第一步的绝对积分总和等于 \(\int_{\mathbb R}|f_b(u)|\,\mathrm du<\infty\)，
第二步用 \(m,n\in\mathbb Z\) 消去因子 \(e^{2\pi\mathrm imn}\)。
这些系数本身快速衰减，所以相应 Fourier 级数绝对一致收敛，定义一个连续函数 \(Q_b\)，且 \(P_b,Q_b\) 的全部 Fourier 系数相同。

为说明“系数相同”确实推出函数相同，可以对差 \(h=P_b-Q_b\) 用核函数
\[
 K_N(x)=\frac1N\left|\sum_{k=0}^{N-1}e^{2\pi\mathrm ikx}\right|^2.
\]
它非负、在 \([0,1]\) 上积分为一；离整数点固定距离时，
等比数列公式给 \(K_N(x)=O(N^{-1})\)。连续周期函数一致连续，把
\(\int_0^1K_N(u)(h(x-u)-h(x))\,\mathrm du\) 分成整数点附近与其余部分，前者由一致连续性控制，后者由上述界趋于零。因此
\(K_N*h\to h\) 一致收敛。另一方面，\(K_N\) 为有限三角多项式，而
\(h\) 的全部 Fourier 系数为零，故 \(K_N*h=0\)，从而 \(h=0\)。
在 \(P_b=Q_b\) 中代入 \(x=a\)，便得到实 \(b\) 时的求和公式。

最后将 \(b\) 延到复数。若 \(b\) 位于一个给定紧集，两侧项的模分别被
\(C e^{-\pi t(n+a)^2+C|n+a|}\) 和
\(C e^{-\pi m^2/t+C|m|}\) 控制。因此两边关于 \(b\) 都整，且在实轴相等。恒等定理给出全部 \(b\in\mathbb C\) 的等式，同时上述估计也证明绝对收敛。
\end{proof}

公式中的两个操作有清楚的对应：空间平移 \(a\) 变成右端的相位
\(e^{2\pi\mathrm ima}\)，空间调制 \(b\) 变成右端频率中心从 \(0\)
移到 \(b\)。取 \(a=1/2,b=0\) 时，左端在半整数格上求和，
右端出现 \((-1)^m\)；这正是 \(\theta_2\) 与 \(\theta_4\) 会在参数反演中交换的原因。

\subsection{模变换}
\label{b3:subsec:240103}

第二十三章的模群由 \(T:\tau\mapsto\tau+1\) 与
\(S:\tau\mapsto-1/\tau\) 生成。对 theta 函数，\(T\) 可以直接在级数中计算，\(S\) 则由刚证明的高斯求和公式产生。空间变量也须随 \(S\) 改为 \(z/\tau\)；只变参数而保持 \(z\) 不动，得不到下面的形式。

\begin{theorem}\label{b3:thm:24-theta-modular}
设 \(z\in\mathbb C\)、\(\tau\in\mathfrak H\)，并定义
\[
 A(\tau)=(-\mathrm i\tau)^{1/2}
   =\exp\left(\tfrac12\operatorname{Log}(-\mathrm i\tau)\right),
 \quad
 B(z,\tau)=A(\tau)e^{\mathrm i z^2/(\pi\tau)}.
\]
这里 \(\operatorname{Log}\) 是右半平面上辐角在
\((-\pi/2,\pi/2)\) 内的全纯对数，因此 \(A(\mathrm it)=\sqrt t>0\)
对所有 \(t>0\) 成立。则
\begin{equation}\label{b3:eq:24-theta-T}
\begin{aligned}
 \theta_1(z,\tau+1)&=e^{\pi\mathrm i/4}\theta_1(z,\tau),&
 \theta_2(z,\tau+1)&=e^{\pi\mathrm i/4}\theta_2(z,\tau),\\
 \theta_3(z,\tau+1)&=\theta_4(z,\tau),&
 \theta_4(z,\tau+1)&=\theta_3(z,\tau),
\end{aligned}
\end{equation}
以及
\begin{equation}\label{b3:eq:24-theta-S}
\begin{aligned}
 \theta_1(z/\tau,-1/\tau)&=-\mathrm iB(z,\tau)\theta_1(z,\tau),\\
 \theta_2(z/\tau,-1/\tau)&=B(z,\tau)\theta_4(z,\tau),\\
 \theta_3(z/\tau,-1/\tau)&=B(z,\tau)\theta_3(z,\tau),\\
 \theta_4(z/\tau,-1/\tau)&=B(z,\tau)\theta_2(z,\tau).
\end{aligned}
\end{equation}
\end{theorem}
\begin{proof}
对整数指标，\(e^{\pi\mathrm in^2}=(-1)^n\)，所以 \(T\) 交换
\(\theta_3,\theta_4\)。对半整数指标，
\[
 e^{\pi\mathrm i(n+1/2)^2}
   =e^{\pi\mathrm i/4}e^{\pi\mathrm in(n+1)}
   =e^{\pi\mathrm i/4},
\]
因为 \(n(n+1)\) 为偶数。这给出 \(\theta_1,\theta_2\) 的相位因子。

为同时算清 \(S\) 的下标和符号，暂记
\[
 \Phi_{a,b}(z,\tau)=
 \sum_{n\in\mathbb Z}e^{\pi\mathrm i\tau(n+a)^2+
                         2\mathrm i(n+a)(z+\pi b)},
 \quad a,b\in\mathbb R.
\]
这只是统一计算的记号；四个函数分别为
\[
 \theta_3=\Phi_{0,0},\quad
 \theta_2=\Phi_{1/2,0},\quad
 \theta_4=\Phi_{0,1/2},\quad
 \theta_1=-\Phi_{1/2,1/2}.
\]
先令 \(\tau=\mathrm it\)、\(t>0\)，在
\cref{b3:lem:24-gaussian-poisson}中把 \(t\) 换成 \(1/t\)，把调制参数
换成 \(b+z/(\pi\mathrm it)\)。该参数可以为复数，引理已经覆盖这种情形。
于是
\[
 \Phi_{a,b}\left(\frac z{\mathrm it},\frac{\mathrm i}t\right)
 =\sqrt t\sum_{m\in\mathbb Z}
    e^{-\pi t(m-b+\mathrm iz/(\pi t))^2}e^{2\pi\mathrm ima}.
\]
展开平方，提出 \(e^{z^2/(\pi t)}\)，再令 \(n=-m\)，得到
\[
 \Phi_{a,b}\left(\frac z{\mathrm it},\frac{\mathrm i}t\right)
 =\sqrt t\,e^{z^2/(\pi t)}e^{2\pi\mathrm iab}
                    \Phi_{b,-a}(z,\mathrm it).
\]
对每个固定 \(z\in\mathbb C\)，等式两边关于 \(\tau\in\mathfrak H\)
都全纯；这一点分别来自正规收敛、\(\tau\ne0\) 和已经选定的
\(A(\tau)\)。在虚轴的全部 \(\tau=\mathrm it\) 上相等，恒等定理便给
\[
 \Phi_{a,b}(z/\tau,-1/\tau)
 =B(z,\tau)e^{2\pi\mathrm iab}\Phi_{b,-a}(z,\tau)
\]
在整个上半平面成立。代入四组 \(a,b\)，即可得到反演公式。
在 \(a=b=1/2\) 的一行还要用
\(\Phi_{1/2,-1/2}=-\Phi_{1/2,1/2}\)，故总因子为
\(e^{\pi\mathrm i/2}(-1)=-\mathrm i\)，而不是 \(\mathrm i\)。
\end{proof}

在 \(\tau=\mathrm it\) 时，反演式中的高斯因子为
\(e^{z^2/(\pi t)}\)。若把等式反解成 \(\theta_3(z,\mathrm it)\)，才出现
衰减因子 \(e^{-z^2/(\pi t)}\)。这个方向区别可以防止记错指数的符号。
整个证明只对实高斯作 Fourier 变换，再作全纯延拓，没有未经估计就平移一条无穷积分直线。
\ProseParagraphBreak
平方根的选取也不能在连续两次变换时随意改变。事实上
\(A(-1/\tau)A(\tau)=1\)：乘积的平方为一，在虚轴取正值，而上半平面连通。
两次 \(S\) 把空间变量 \(z\) 变为 \(-z\)，高斯因子互相抵消。
三个偶函数的结果不变，\(\theta_1\) 的两个因子
\((-\mathrm i)^2=-1\) 则正好恢复它的奇性。这为根分支和相位提供了一个直接核对。

\subsection{\texorpdfstring{\(\theta\)}{\AnalysisPdfTheta} 常数}
\label{b3:subsec:240104}

\begin{definition}\label{b3:def:24-theta-constants}
对 \(\tau\in\mathfrak H\)，四个值
\(\theta_j(0,\tau)\)、\(1\le j\le4\)，称为\term[theta-changshu]{theta 常数}[theta constant]。
当空间变量已固定为零且参数明确时，简记为
\(\theta_j(\tau)\) 或 \(\theta_j\)。“常数”只表示不再依赖 \(z\)，并不表示它们不随 \(\tau\) 改变。
\end{definition}

其中 \(\theta_1(\tau)=0\)。另外三个常数的展开为
\begin{equation}\label{b3:eq:24-theta-constant-series}
\begin{aligned}
 \theta_2(\tau)&=2q^{1/4}(1+q^2+q^6+q^{12}+\cdots),\\
 \theta_3(\tau)&=1+2q+2q^4+2q^9+\cdots,\\
 \theta_4(\tau)&=1-2q+2q^4-2q^9+\cdots,
 \quad q=e^{\pi\mathrm i\tau}.
\end{aligned}
\end{equation}
第一行括号中的一般指数为 \(n(n+1)\)、\(n\ge0\)，后两行的指数为
\(n^2\)、\(n\ge1\)。这三者在上半平面都不为零；该事实将在下一节由乘积公式证明，目前尚不凭展开的首项排除其余项的复数抵消。
\ProseParagraphBreak
在\cref{b3:thm:24-theta-modular}中取 \(z=0\)，得到常数反演
\[
 \theta_2(-1/\tau)=A(\tau)\theta_4(\tau),\quad
 \theta_3(-1/\tau)=A(\tau)\theta_3(\tau),\quad
 \theta_4(-1/\tau)=A(\tau)\theta_2(\tau).
\]
这解释了为何 theta 常数与模形式有关，却不能直接把每个常数称作
第二十三章意义下的全模群标量模形式：变换因子出现平方根，而且
\(T\) 与 \(S\) 会交换不同函数。尤其 \(\theta_3(\tau+1)=\theta_4(\tau)\)，
二者的首个非恒定系数分别为 \(2q\) 和 \(-2q\)，所以并不相等。
把函数取幂并作适当组合，才可能消除这些相位和交换；只凭它含有一个
\((-\mathrm i\tau)^{1/2}\) 因子，不能省去其余变换条件。
\ProseParagraphBreak
分数幂的单值性同样可从 \(T\) 看清。虽然 \(\tau\mapsto\tau+2\) 保持
\(q=e^{\pi\mathrm i\tau}\) 不变，却使
\(\theta_2(\tau+2)=\mathrm i\theta_2(\tau)\)，因为
\(q^{1/4}=e^{\pi\mathrm i\tau/4}\) 随参数路径改变。后文的
\(q\partial_q\) 是已在\cref{b3:prop:24-theta-heat}说明的局部微分算子，
不会把这个全局分支信息丢掉。
\ProseParagraphBreak
参数反演还有一个可计算的用途。设 \(t>0\)，由于
\(n^2\ge1+3(n-1)\) 对 \(n\ge1\) 成立，
\[
 0<\theta_3(\mathrm it)-1
       =2\sum_{n\ge1}e^{-\pi tn^2}
       \le\frac{2e^{-\pi t}}{1-e^{-3\pi t}}.
\]
大 \(t\) 时原级数很快收敛，小 \(t\) 时则用
\[
 \theta_3(\mathrm it)=t^{-1/2}\theta_3(\mathrm i/t)
     =t^{-1/2}\left(1+O(e^{-\pi/t})\right)
     \quad(t\downarrow0).
\]
余项直接由刚给出的界把 \(t\) 换成 \(1/t\) 得到。这不只是一个形式对称：
它把收敛慢的参数变成收敛快的参数，同时精确给出主要大小
\(t^{-1/2}\)。从平方和生成函数走向 Mellin 变换时，零端点的这个大小也将决定积分最初的收敛范围。

% !TeX root = ../../Book3.tex
\section{Jacobi 三重积公式}
\label{b3:sec:2402}

级数把不同整数指标的贡献相加，适合研究变换和表示数；乘积把零点逐个列出，适合判断除法是否有意义。这里要把上一节的 theta 级数改写成乘积，并说明为什么无穷多个非零因子不会在极限中另外制造零点。

证明从有限多项式开始。先把有限乘积的系数算清楚，再把指标居中，使指数变成平方，最后在紧集上取极限。这条路线同时给出恒等式中的常数，避免只凭零点相同就猜测两个函数相等。

\subsection{三重积恒等式}
\label{b3:subsec:240201}

下面的有限恒等式起着普通二项式公式的作用。记号中的下标 \(Q\) 表示因子按等比数列变化；这里的系数一般是复数，不借用计数或正性来证明公式。

\Needspace{8\baselineskip}
\begin{lemma}[有限二项式恒等式]\label{b3:lem:24-finite-binomial}
设 \(Q\in\mathbb C\)、\(0<|Q|<1\)，并令
\[
 (Q;Q)_0=1,\qquad (Q;Q)_N=\prod_{j=1}^N(1-Q^j)
 \quad(N\ge1).
\]
对整数 \(N\ge0\) 及 \(0\le r\le N\)，定义
\[
 {N\brack r}_Q
 =\frac{(Q;Q)_N}{(Q;Q)_r(Q;Q)_{N-r}},
\]
并在 \(r<0\) 或 \(r>N\) 时规定 \({N\brack r}_Q=0\)。则对任意 \(t\in\mathbb C\)，
\begin{equation}\label{b3:eq:24-finite-binomial}
 \prod_{j=0}^{N-1}(1+tQ^j)
 =\sum_{r=0}^N Q^{r(r-1)/2}{N\brack r}_Q\,t^r.
\end{equation}
当 \(N=0\) 时，左边的空乘积和右边都等于一。
\end{lemma}
\begin{proof}
因 \(|Q|<1\)，所有分母因子都不为零。先证明递推关系
\[
 {N+1\brack r}_Q
 ={N\brack r}_Q+Q^{N+1-r}{N\brack r-1}_Q.
\]
在 \(1\le r\le N\) 时，把右边通分为分母
\((Q;Q)_r(Q;Q)_{N+1-r}\)，分子为
\[
 (Q;Q)_N\bigl((1-Q^{N+1-r})
              +Q^{N+1-r}(1-Q^r)\bigr)
 =(Q;Q)_{N+1}.
\]
在 \(r=0,N+1\) 时，越界系数为零的约定直接给出同一递推。

现在对 \(N\) 归纳。把 \(N\) 个因子的展开乘以 \(1+tQ^N\)，所得 \(t^r\) 的系数为
\[
 Q^{r(r-1)/2}{N\brack r}_Q
 +Q^N Q^{(r-1)(r-2)/2}{N\brack r-1}_Q.
\]
提出 \(Q^{r(r-1)/2}\) 后，括号中的式子就是上述递推的右边，因此得到 \(N+1\) 的公式。
\end{proof}

例如 \(N=2\) 时，公式是
\[
 (1+t)(1+tQ)=1+(1+Q)t+Qt^2.
\]
额外的因子 \(Q^{r(r-1)/2}\) 不能省去：选择两个因子中的 \(t\) 项时，所得系数是 \(Q\)，而不是一。

将有限公式居中后，就得到下面的\term[jacobi-sanchongji]{Jacobi 三重积公式}[Jacobi triple product identity]。

\Needspace{8\baselineskip}
\begin{theorem}[Jacobi 三重积]\label{b3:thm:24-triple-product}
设 \(q,w\in\mathbb C\)，且 \(0<|q|<1\)、\(w\ne0\)。则
\[
 \sum_{n\in\mathbb Z}q^{n^2}w^n
 =\prod_{m=1}^{\infty}
   (1-q^{2m})(1+wq^{2m-1})(1+w^{-1}q^{2m-1}).
\]
固定 \(q\) 后，级数绝对正规收敛，乘积正规收敛，二者都在
\(\mathbb C\setminus\{0\}\) 上定义关于 \(w\) 的全纯函数。
\end{theorem}
\begin{proof}
在式\eqref{b3:eq:24-finite-binomial}中，把因子个数取为 \(2N\)，再取
\[
 Q=q^2,\qquad t=wq^{-(2N-1)}.
\]
左边各因子中 \(q\) 的指数依次为
\(-(2N-1),-(2N-3),\ldots,2N-1\)。对负指数的那一半使用
\[
 1+wq^{-(2j-1)}
 =wq^{-(2j-1)}(1+w^{-1}q^{2j-1}).
\]
提出的全部单项式相乘为 \(w^Nq^{-N^2}\)，因为前 \(N\) 个正奇数的和为 \(N^2\)。

右边第 \(\ell\) 项含有
\[
 q^{\ell(\ell-1)}w^\ell q^{-(2N-1)\ell}
 =w^\ell q^{\ell^2-2N\ell}.
\]
因此两边乘以 \(w^{-N}q^{N^2}\)，并令 \(r=\ell-N\)，就得到
\begin{equation}\label{b3:eq:24-centered-product}
 \prod_{j=1}^N(1+wq^{2j-1})(1+w^{-1}q^{2j-1})
 =\sum_{r=-N}^N q^{r^2}w^r{2N\brack N+r}_{q^2}.
\end{equation}
平方指数来自 \(\ell^2-2N\ell+N^2=(\ell-N)^2\)，而双边指标来自把长度为 \(2N\) 的有限区间移到原点两侧。

接着说明极限为什么可以进入右边的和。由
\(\sum_{j\ge1}|q|^{2j}<\infty\) 和\cref{b3:prop:20-absolute-product}的非零乘积判据，
\[
 A=\prod_{j=1}^{\infty}(1-q^{2j})
\]
存在且非零。对每个固定整数 \(r\)，当 \(N\to\infty\) 时，商式中的三个有限乘积都趋于 \(A\)，故
\[
 {2N\brack N+r}_{q^2}\longrightarrow A^{-1}.
\]
这里还需要一个不随 \(N,r\) 变化的界。记
\[
 d=\prod_{j=1}^{\infty}(1-|q|^{2j})>0,\qquad
 M=\prod_{j=1}^{\infty}(1+|q|^{2j})<\infty.
\]
任一分母有限乘积的模不小于 \(d\)，分子有限乘积的模不大于 \(M\)，所以
\[
 \left|{2N\brack N+r}_{q^2}\right|
 \le Md^{-2}\qquad(-N\le r\le N).
\]

固定紧圆环 \(a\le|w|\le b\)，其中 \(0<a\le b<\infty\)，令
\(B=\max\{b,a^{-1},1\}\)。把区间外的系数补为零后，右边每一项的模都不超过
\[
 Md^{-2}|q|^{r^2}B^{|r|},\qquad r\in\mathbb Z.
\]
这个双边数列可和，因为平方衰减压过固定的指数增长。于是先把尾和统一压小，再对剩余有限项取极限，得到右边在紧圆环上一致趋于
\(A^{-1}\sum_{r\in\mathbb Z}q^{r^2}w^r\)。

左边也在紧圆环上一致收敛：两个因子与一的差分别被
\(b|q|^{2j-1}\) 和 \(a^{-1}|q|^{2j-1}\) 控制，控制数列均可和。
\cref{b3:prop:20-normal-product}因而适用。对式\eqref{b3:eq:24-centered-product}取极限，再乘以 \(A\)，即得所需公式及收敛结论。
\end{proof}

固定 \(r\) 的系数极限并不足以独自证明定理，因为求和区间还在扩大。证明中的共同上界专门控制那些随 \(N\) 向两端移动的项。另一方面，三重积的某个因子可以为零；正规收敛所保证的是充分远的尾部非零，并没有排除有限头部给出的零点。

\subsection{乘积展开}
\label{b3:subsec:240202}

现在恢复 \(\tau\in\mathfrak H\)、\(q=e^{\pi\mathrm i\tau}\)。本章的空间变量是 \(z\)，因此周期中出现 \(\pi\)，而不是一。以下所有
\(q^{1/4}\) 都表示 \(e^{\pi\mathrm i\tau/4}\)；把 \(q\) 单独看成穿孔圆盘中的变量时，不能另选一个全局的四次根来替换它。

\Needspace{8\baselineskip}
\begin{proposition}[四种 theta 函数的乘积]\label{b3:prop:24-theta-products}
设 \(\tau\in\mathfrak H\)、\(q=e^{\pi\mathrm i\tau}\)，并以
\(q^{1/4}=e^{\pi\mathrm i\tau/4}\) 规定分数幂。则上一节定义的四种 theta 函数对每个 \(z\in\mathbb C\) 满足
\begin{align}
 \theta_1(z,\tau)
 &=2q^{1/4}\sin z
   \prod_{n\ge1}(1-q^{2n})
   (1-2q^{2n}\cos2z+q^{4n}),
   \label{b3:eq:24-theta1-product}\\
 \theta_2(z,\tau)
 &=2q^{1/4}\cos z
   \prod_{n\ge1}(1-q^{2n})
   (1+2q^{2n}\cos2z+q^{4n}),
   \label{b3:eq:24-theta2-product}\\
 \theta_3(z,\tau)
 &=\prod_{n\ge1}(1-q^{2n})
   (1+2q^{2n-1}\cos2z+q^{4n-2}),
   \label{b3:eq:24-theta3-product}\\
 \theta_4(z,\tau)
 &=\prod_{n\ge1}(1-q^{2n})
   (1-2q^{2n-1}\cos2z+q^{4n-2}).
   \label{b3:eq:24-theta4-product}
\end{align}
固定 \(\tau\) 后，乘积关于 \(z\) 正规收敛，并可对 \(z\) 作任意固定阶数的逐项微分。其收敛在 \((z,\tau)\) 的紧集上也一致。
\end{proposition}
\begin{proof}
在三重积中分别取 \(w=e^{2\mathrm iz}\) 和 \(w=-e^{2\mathrm iz}\)，得到 \(\theta_3,\theta_4\) 的式子。其中两个含 \(w\) 的因子相乘，分别给出正号和负号的余弦项。

对半整数指标，先在级数中提取共同因子：
\[
 \begin{aligned}
 \theta_2(z,\tau)
 &=q^{1/4}e^{\mathrm iz}
   \sum_{n\in\mathbb Z}q^{n^2}(qe^{2\mathrm iz})^n,\\
 \theta_1(z,\tau)
 &=-\mathrm i q^{1/4}e^{\mathrm iz}
   \sum_{n\in\mathbb Z}q^{n^2}(-qe^{2\mathrm iz})^n.
 \end{aligned}
\]
第一行的两个指数因子成为
\(1+q^{2m}e^{2\mathrm iz}\) 和
\(1+q^{2m-2}e^{-2\mathrm iz}\)。
后一个乘积在 \(m=1\) 时提供 \(1+e^{-2\mathrm iz}\)；将剩余指标减一，便与第一个乘积配对。外面的因子化为
\[
 e^{\mathrm iz}(1+e^{-2\mathrm iz})=2\cos z.
\]
第二行同样给出负号的两个指数因子，而外面的因子为
\[
 -\mathrm i e^{\mathrm iz}(1-e^{-2\mathrm iz})=2\sin z.
\]
这就解释了四个公式中不同的首因子及符号。

若 \((z,\tau)\) 限制在一个紧集，则有 \(|q|\le r<1\) 及
\(|e^{\pm2\mathrm iz}|\le C\)。把每个二次因子拆为
\(1\pm q^m e^{2\mathrm iz}\) 与 \(1\pm q^m e^{-2\mathrm iz}\)，其偏离一的量被 \(Cr^m\) 控制。故乘积的正规收敛及有限分组均由\cref{b3:prop:20-normal-product}保证。固定 \(\tau\) 后，在稍大的 \(z\) 紧集上使用同一估计及 Cauchy 公式，就得到各阶导数的局部一致收敛。
\end{proof}

乘积表示给出的不只是一些显然存在的零点，还包括所有零点及其重数。后文把 theta 函数放在分母时，将用到这一完整信息。

\Needspace{8\baselineskip}
\begin{proposition}[theta 零点及常数非零性]\label{b3:prop:24-theta-zeros}
设 \(\tau\in\mathfrak H\)，令
\(\Lambda=\pi\mathbb Z+\pi\tau\mathbb Z\)。四个整函数
\(z\mapsto\theta_j(z,\tau)\) 的零点集分别为
\[
 \begin{aligned}
 Z(\theta_1)&=\Lambda,
 &Z(\theta_2)&=\frac{\pi}{2}+\Lambda,\\
 Z(\theta_3)&=\frac{\pi+\pi\tau}{2}+\Lambda,
 &Z(\theta_4)&=\frac{\pi\tau}{2}+\Lambda.
 \end{aligned}
\]
其中每一个零点都是单零点。特别地，
\(\theta_2(0,\tau)\)、\(\theta_3(0,\tau)\)、\(\theta_4(0,\tau)\)
均不为零，而 \(\theta_1(0,\tau)=0\) 且
\(\partial_z\theta_1(0,\tau)\ne0\)。
\end{proposition}
\begin{proof}
在 \(\theta_1\) 的乘积中，因子 \(1-q^{2n}\) 从不为零。
首因子 \(\sin z\) 的零点是 \(\pi\mathbb Z\)，且都单重。另外两个因子为
\[
 1-q^{2n}e^{2\mathrm iz},\qquad
 1-q^{2n}e^{-2\mathrm iz}.
\]
它们的零点分别是 \(-n\pi\tau+\pi\mathbb Z\) 和
\(n\pi\tau+\pi\mathbb Z\)。在各自的零点处，导数为
\(-2\mathrm i\) 或 \(2\mathrm i\)，所以也都是单零点。
因 \(\operatorname{Im}\tau>0\)，这些水平直线彼此不同，且不与
\(\sin z\) 的零点重叠。于是每个格点恰好来自一个单零因子。

对 \(\theta_2\)，首因子换成 \(\cos z\)，两个指数因子换成正号，零点的实部相应平移 \(\pi/2\)，得到 \(\pi/2+\Lambda\)。
对 \(\theta_3,\theta_4\)，指数为 \(2n-1\)，所以虚方向的平移为
\(\pm(n-\tfrac12)\pi\tau\)。正号因子再增加 \(\pi/2\) 的水平平移，负号因子不增加，分别得到所列的两个集合。各因子的导数仍非零，不同因子的零点仍不重叠。

最后，正规乘积在每个小邻域都可写成有限头部乘一个非零的全纯尾部；这是\cref{b3:prop:20-product-zero-orders}的具体适用情形。因此没有遗漏的额外零点，也没有因无限乘积而增加的重数。原点不属于后三个格的平移：比较虚部先确定 \(\tau\) 的系数，再比较实部即可验证。常数非零性随之成立。
\end{proof}

四个零点集在一个格胞中分别占据四个半周期位置。准周期平移会改变函数值的非零乘子，但不会改变零点的重数；因此这里描述的是整张平面上的零点排列，而不是只在一个基本区域内的局部结论。

\subsection{特殊值}
\label{b3:subsec:240203}

令 \(z=0\)。为减轻记号，以下把 \(\theta_j(0,\tau)\) 简写成
\(\theta_j\)，但撇号始终表示对空间变量 \(z\) 求导。由乘积公式，
\[
 \begin{aligned}
 \theta_2&=2q^{1/4}
       \prod_{n\ge1}(1-q^{2n})(1+q^{2n})^2,\\
 \theta_3&=\prod_{n\ge1}(1-q^{2n})(1+q^{2n-1})^2,\\
 \theta_4&=\prod_{n\ge1}(1-q^{2n})(1-q^{2n-1})^2.
 \end{aligned}
\]
这些是关于参数 \(\tau\) 的函数；“常数”只表示空间变量已经取零。

三个非零常数还决定奇函数在原点的导数，这一关系称为\term[jacobi-daoshu-hengdengshi]{Jacobi 导数恒等式}[Jacobi derivative identity]。

\Needspace{6\baselineskip}
\begin{proposition}[Jacobi 导数恒等式]\label{b3:prop:24-jacobi-derivative}
设 \(\tau\in\mathfrak H\)，四个 theta 函数采用本章的级数归一化，则
\begin{equation}\label{b3:eq:24-derivative-identity}
 \partial_z\theta_1(0,\tau)
 =\theta_2(0,\tau)\theta_3(0,\tau)\theta_4(0,\tau).
\end{equation}
两边均不为零。
\end{proposition}
\begin{proof}
令 \(q=e^{\pi\mathrm i\tau}\)，并规定
\(q^{1/4}=e^{\pi\mathrm i\tau/4}\)。在式\eqref{b3:eq:24-theta1-product}中除以 \(z\) 并令 \(z\to0\)，由于乘积局部一致收敛及 \(\sin z/z\to1\)，得到
\[
 \theta_1'(0,\tau)
 =2q^{1/4}\prod_{n\ge1}(1-q^{2n})^3.
\]
现在将三个 theta 常数的乘积式相乘。共同因子已经给出右边这个式子，剩下的因子是
\[
 \left(\prod_{n\ge1}
    (1+q^{2n})(1-q^{4n-2})\right)^2.
\]
记 \(A=\prod_{n\ge1}(1-q^{2n})\ne0\)。由
\(1+q^{2n}=(1-q^{4n})/(1-q^{2n})\)，括号内的乘积等于
\[
 \frac{\displaystyle
        \prod_{n\ge1}(1-q^{4n})
        \prod_{n\ge1}(1-q^{4n-2})}
      {\displaystyle\prod_{n\ge1}(1-q^{2n})}=1.
\]
分子两列指数 \(4,8,12,\ldots\) 与 \(2,6,10,\ldots\) 正好分割全部正偶数。这里的分组和约分合法，因为各乘积绝对收敛且非零；并没有把两个发散乘积相除。恒等式随之成立，非零性也从这些乘积直接得到。
\end{proof}

这个归一化没有额外的 \(\pi\) 因子，因为本章级数含
\(e^{2\mathrm inz}\)，基本水平准周期为 \(\pi\)。若改用周期为一的空间坐标，空间求导的换元因子也必须随之改变。

当 \(\tau=\mathrm it\)、\(t>0\) 时，\(q=e^{-\pi t}\in(0,1)\)。
三个常数乘积中的全部因子都是正数，因此
\(\theta_2,\theta_3,\theta_4>0\)。尤其是交替级数
\[
 \theta_4=1-2q+2q^4-2q^9+\cdots
\]
始终为正。单看前两项，在 \(q>1/2\) 时容易作出错误判断；乘积同时控制了被截去的无限尾部。对一般复参数，仍有非零性，但不再谈正负。

在 \(q\to0\) 时，级数还给出
\[
 \begin{aligned}
 \theta_2&=2q^{1/4}\bigl(1+q^2+O(q^6)\bigr),\\
 \theta_3&=1+2q+O(q^4),\qquad
 \theta_4=1-2q+O(q^4).
 \end{aligned}
\]
因此它们的乘积为 \(2q^{1/4}(1-3q^2+O(q^4))\)，与导数乘积的展开相符。这个检查能辨认归一化和低次系数，却不能代替对所有 \(\tau\in\mathfrak H\) 成立的恒等式证明。

\subsection{三重积的系数、零点与对数导数}
\label{b3:subsec:240204}

三重积允许比较系数，但要先说清楚比较的是哪个变量。固定
\(w\ne0\) 后，三重积两边关于 \(q\) 在 \(|q|<1\) 全纯，并在
\(q=0\) 都取值一；局部一致收敛保证可以使用幂级数的唯一性。
固定 \(q\) 后，关于 \(w\) 则是穿孔平面上的 Laurent 展开。两种比较都依赖解析收敛。靠近 \(|q|=1\) 时，上面的共同控制不再一致，不能未经新的估计就把 \(q\) 代成单位根。

还有一种需要格外区分的操作：对乘积求对数导数。含有零点的整函数没有全平面上的全纯对数，但在每个不含零点的小圆盘中都有局部对数；不同分支只差常数，所以对数的导数不依赖分支。

\Needspace{8\baselineskip}
\begin{proposition}[theta 的对数导数]\label{b3:prop:24-theta-log-derivative}
设 \(\tau\in\mathfrak H\)、\(q=e^{\pi\mathrm i\tau}\)，令
\(\Lambda=\pi\mathbb Z+\pi\tau\mathbb Z\) 及
\[
 L(z)=\frac{\partial_z\theta_1(z,\tau)}{\theta_1(z,\tau)}
 \qquad(z\in\mathbb C\setminus\Lambda).
\]
则
\begin{equation}\label{b3:eq:24-theta-logarithmic-series}
 L(z)=\cot z+
 4\sin(2z)\sum_{n\ge1}
 \frac{q^{2n}}{1-2q^{2n}\cos(2z)+q^{4n}}.
\end{equation}
右边的级数在 \(\mathbb C\setminus\Lambda\) 的紧集上绝对一致收敛，并可作任意固定阶数的逐项微分。在更小的范围
\(|\operatorname{Im}z|<\pi\operatorname{Im}\tau\)、\(z\notin\pi\mathbb Z\) 内，还有
\begin{equation}\label{b3:eq:24-theta-log-fourier}
 L(z)=\cot z+
 4\sum_{k\ge1}\frac{q^{2k}}{1-q^{2k}}\sin(2kz).
\end{equation}
函数 \(-L'\) 是以 \(\pi,\pi\tau\) 为周期的亚纯函数，其全部极点是
\(\Lambda\)，在每个 \(\lambda\in\Lambda\) 处的主部都是
\((z-\lambda)^{-2}\)。
\end{proposition}
\begin{proof}
在避开 \(\Lambda\) 的紧集上，theta 乘积的有限头部非零，尾因子一致接近一。把紧集稍微放大后，尾因子的偏差仍由几何级数控制。因此\cref{b3:prop:20-product-log-derivative}的条件成立：可对局部对数求导，并可继续逐项微分。由
\[
 \frac{\mathrm d}{\mathrm dz}
   \log(1-2q^{2n}\cos2z+q^{4n})
 =\frac{4q^{2n}\sin2z}
        {1-2q^{2n}\cos2z+q^{4n}},
\]
以及首因子 \(\sin z\) 的对数导数，得到第一式。这里的公式只要求避开零点，不要求 \(z\) 留在一个水平带内。

第二式需要额外的几何级数展开。若
\(|\operatorname{Im}z|<-\log|q|=\pi\operatorname{Im}\tau\)，则对每个 \(n\ge1\)，
\(|q^{2n}e^{\pm2\mathrm iz}|<1\)。将两个指数因子的对数导数分别展开后相加，得到
\[
 \frac{-2\mathrm i q^{2n}e^{2\mathrm iz}}
           {1-q^{2n}e^{2\mathrm iz}}
 +\frac{2\mathrm i q^{2n}e^{-2\mathrm iz}}
           {1-q^{2n}e^{-2\mathrm iz}}
 =4\sum_{k\ge1}q^{2nk}\sin(2kz).
\]
在闭子带 \(|\operatorname{Im}z|\le R<-\log|q|\) 上，相应双重级数的绝对值和不超过常数乘以
\[
 \sum_{k\ge1}
   \frac{|q|^{2k}e^{2kR}}{1-|q|^{2k}}<\infty.
\]
故可交换 \(n,k\) 的求和顺序，先对 \(n\) 求和，就得到式\eqref{b3:eq:24-theta-log-fourier}。

由上一节的准周期公式，
\[
 \theta_1(z+\pi,\tau)=-\theta_1(z,\tau),\qquad
 \theta_1(z+\pi\tau,\tau)
 =-q^{-1}e^{-2\mathrm iz}\theta_1(z,\tau).
\]
取对数导数即得 \(L(z+\pi)=L(z)\) 和
\(L(z+\pi\tau)=L(z)-2\mathrm i\)，再求导便消去加法常数。
另一方面，\cref{b3:prop:24-theta-zeros}保证每个格点均为单零点。
在 \(\lambda\) 附近可写
\(\theta_1(z,\tau)=(z-\lambda)h(z)\)，其中 \(h\) 全纯且非零。于是
\[
 L(z)=\frac1{z-\lambda}+\frac{h'(z)}{h(z)},\qquad
 -L'(z)=\frac1{(z-\lambda)^2}
          -\left(\frac{h'}h\right)'(z).
\]
第二项全纯，从而得到全部极点和主部。
\end{proof}

例如，在原点附近对式\eqref{b3:eq:24-theta-log-fourier}求导，并使用
\(\csc^2z=z^{-2}+\tfrac13+O(z^2)\)，得到
\begin{equation}\label{b3:eq:24-theta-local-part}
 -L'(z)=\frac1{z^2}+\frac13
       -8\sum_{k\ge1}\frac{kq^{2k}}{1-q^{2k}}+O(z^2).
\end{equation}
这里极点的系数一已经由单零点确定，常数项则由乘积归一化决定。后文与同一格的 \(\wp\) 函数比较时，二者之差只能剩下常数；不能只因二者都双周期就认定它们相等。

参数方向的对数微分也要有共同控制。例如
\[
 B(q)=\prod_{n\ge1}(1-q^{2n})\qquad(|q|<1)
\]
处处非零且全纯，并满足
\[
 q\frac{B'(q)}{B(q)}
 =-\sum_{n\ge1}\frac{2nq^{2n}}{1-q^{2n}}.
\]
在 \(|q|\le r<1\) 上，右边由
\(\sum_{n\ge1}2nr^{2n}/(1-r^{2n})<\infty\) 控制。
其他 theta 常数乘积的参数导数有同样的几何衰减。若出现
\(q^{1/4}\)，则在 \(\tau\) 平面使用
\[
 q\partial_q=\frac1{\pi\mathrm i}\partial_\tau,
 \qquad q\partial_q\log(q^{1/4})=\frac14.
\]
后一式是局部对数的导数陈述，不要求在整个穿孔 \(q\) 圆盘选定单值的 \(\log q\)。这些估计将使后文从乘积转到约数和的逐项运算有明确的适用范围。

% !TeX root = ../../Book3.tex
\section{\texorpdfstring{\(\theta\)}{\AnalysisPdfTheta} 级数作为生成函数}
\label{b3:sec:2403}

系数计数必须先约定是否区分顺序和符号。theta 乘法自动区分各坐标，而且正负整数分别出现，因此它计算的是有序带符号表示数。

\subsection{生成函数}
\label{b3:subsec:240301}

设 \(r_k(n)=\#\{(a_1,\ldots,a_k)\in\mathbb Z^k:\sum a_j^2=n\}\)。每个坐标绝对值不超过 \(\sqrt n\)，故数目有限。对 \(|q|<1\)，绝对收敛保证
\[
\theta_3(0,\tau)^k=\sum_{n\ge0}r_k(n)q^n.
\]
此处 \(r_k(0)=1\)。例如 \(r_2(1)=4\)、\(r_4(1)=8\)，分别选择非零坐标的位置和它的符号，不能把它们解释为不计顺序的分拆数。

\subsection{乘积与卷积}
\label{b3:subsec:240302}

乘法给 \(r_{k+\ell}(n)=\sum_{m=0}^n r_k(m)r_\ell(n-m)\)。这是把前 \(k\) 个坐标与后 \(\ell\) 个坐标分别计数的精确分解。

\ProseParagraphBreak

若改用 \(\theta_4\)，每个坐标带符号 \((-1)^{a_j}\)。由于 \(a_j^2\equiv a_j\pmod2\)，同一层 \(\sum a_j^2=n\) 上总符号恒为 \((-1)^n\)。因此 \(\theta_4^k=\sum(-1)^nr_k(n)q^n\)。它筛选奇偶层，却不单独筛选每个坐标的奇偶模式。

\subsection{系数比较}
\label{b3:subsec:240303}

在四个变量中还有一个有用的精确变换。对 \(a=(a_1,a_2,a_3,a_4)\) 用正交矩阵
\[
\frac12\begin{pmatrix}1&1&1&1\\1&1&-1&-1\\1&-1&1&-1\\1&-1&-1&1\end{pmatrix}
\]
变换为 \(b\)。当 \(\sum a_j\) 为奇数时，所有 \(b_j\in\mathbb Z+1/2\)；反向变换相同。半整数四元组又分为两类：写 \(b_j=m_j+1/2\)，按 \(\sum m_j\) 的奇偶分。只有其中一类在逆变换下给整数向量；将一个 \(b_j\) 改号交换两类且保持平方和。故全部半整数向量的计数恰为奇数平方和整数向量计数的两倍。比较级数得到
\begin{equation}\label{b3:eq:24-quartic-identity}
\theta_3^4=\theta_4^4+\theta_2^4.
\end{equation}
这个推导既说明常数二从何而来，也说明为什么四次方出现。它不是把前几项相同当成无穷恒等式，而是给出保持范数的对应。

\ProseParagraphBreak

生成函数还可以筛选同余条件。若 \(A(q)=\sum a_nq^n\)，\(\zeta_m=e^{2\pi\mathrm{i}/m}\)，则有限 Fourier 正交关系给
\[
 \frac1m\sum_{j=0}^{m-1}\zeta_m^{-bj}A(\zeta_m^j q)
 =\sum_{n\equiv b\ (m)}a_nq^n.
\]
交换有限和与级数不需要新的收敛定理；只要 \(|q|<1\)，每个旋转参数仍在原收敛圆盘内。对平方和而言，这个操作筛选总平方和的剩余类。若要每个坐标分别满足某些同余条件，则须先对每个一维 theta 因子筛选，再相乘。两个程序统计的集合不同，不能仅凭最后都含若干符号因子就互换。
\ProseParagraphBreak
四维正交变换中，半整数层的两类也可直接检查。写 \(b_j=m_j+1/2\)，逆变换的第一坐标为 \((\sum m_j+2)/2\)，其余坐标具有同样的整数奇偶条件。因此恰在 \(\sum m_j\) 偶数时，逆像为整数四元组；这时逆像的坐标和为 \(2b_1\)，必为奇数。把一个 \(b_j\) 改号，相应的整数 \(m_j\) 变成 \(-m_j-1\)，总和奇偶翻转，所以两类等势且逐层保范数。这样便把前述系数比较中的二倍关系落实为显式双射，而非只比较总的无穷集合大小。
\ProseParagraphBreak
平均表示数可以从几何体积作一个独立的数量级检验。令 \(A_k(R)\) 为 \(\mathbb Z^k\) 中长度不超过 \(R\) 的点数。以每个格点为中心放置边长一的闭立方体，它们内部不交，每个体积一。所取立方体的并包含半径 \(R-\sqrt{k}/2\) 的球，又包含在半径 \(R+\sqrt{k}/2\) 的球内。因此
\[
 A_k(R)=V_kR^k+O_k(R^{k-1}),
\]
其中 \(V_k\) 为单位球体积。这个估计不控制单个球壳上的点数，却控制从0累计到 \(R^2\) 的表示数。令 \(X=R^2\)，即为 \(\sum_{n\le X}r_k(n)=V_kX^{k/2}+O_k(X^{(k-1)/2})\)。平方和问题中，逐项精确公式与累计体积估计是不同层次的结果。

% !TeX root = ../../Book3.tex
\section{二平方和与四平方和}
\label{b3:sec:2404}

二平方和可以从 Gaussian 整数的唯一分解得到；四平方和的精确计数则可利用 theta 乘积的对数导数。后者需要一个半周期差值恒等式，下面将它从前章的椭圆函数推出来。

\subsection{二平方和定理}
\label{b3:subsec:240401}

\begin{theorem}[二平方和]\label{b3:thm:24-two-squares}
正整数 \(n\) 可表示为两个整数平方之和，当且仅当每个 \(p\equiv3\pmod4\) 的素因子在 \(n\) 中的指数为偶数。计数公式为
\(r_2(n)=4\sum_{d\mid n}\chi_4(d)\)，其中 \(\chi_4(d)=0\) 对偶数，\(\chi_4(d)=(-1)^{(d-1)/2}\) 对奇数。
\end{theorem}
\begin{proof}
先说明所需代数。Gaussian 整数 \(\mathbb Z[\mathrm{i}]\) 的范数为 \(a^2+b^2\)。对商的实虚部各取最近整数，可使除法余项的范数至多为除数范数的一半，因此存在 Euclid 算法、最大公因子和 Bézout 恒等式。若不可约元整除积，Bézout 恒等式证明它整除某个因子，继而按范数归纳得到唯一分解。
有理素数 \(p\equiv3\pmod4\) 仍不可约：若分解，其因子范数只能为 \(p\)，但平方模四不可能相加得三。对 \(p\equiv1\pmod4\)，Wilson 配对证明 \((p-1)!\equiv-1\pmod p\)，又有 \((p-1)!\equiv(-1)^{(p-1)/2}(((p-1)/2)!)^2\)，所以存在 \(u^2\equiv-1\pmod p\)。若 \(p\) 在 Gaussian 整数中不可约，它整除 \((u+\mathrm{i})(u-\mathrm{i})\) 就会整除一因子，这在虚部为 \(\pm1\) 时不可能。于是 \(p=\pi\overline\pi\)，其中 \(N\pi=p\)，两因子不相伴。素数二则为 \(-\mathrm{i}(1+\mathrm{i})^2\)。
设 \(n=2^a\prod_{p\equiv1(4)}p^{a_p}\prod_{q\equiv3(4)}q^{b_q}\)。若 \(N\alpha=n\)，每个惰性素数 \(q\) 在范数中的指数必须偶；反之这些指数偶时，可以给 \(\alpha\) 选 \(q^{b_q/2}\)、\((1+\mathrm{i})^a\)，并对每个分裂素数在 \(\pi\) 与 \(\overline\pi\) 间分配共 \(a_p\) 个因子。因此表示数为 \(4\prod(a_p+1)\)，前面的四为四个单位。约数和 \(\sum_{d\mid n}\chi_4(d)\) 按素数幂相乘：二贡献一，分裂素数贡献 \(a_p+1\)，惰性素数贡献一或零，恰好得到公式。
\end{proof}

\subsection{Jacobi 四平方和定理}
\label{b3:subsec:240402}

\begin{theorem}[Jacobi 四平方和]\label{b3:thm:24-four-squares}
对每个正整数 \(n\)，有
\[
r_4(n)=8\sum_{d\mid n,\ 4\nmid d}d.
\]
特别地，每个正整数均为四个整数平方之和。
\end{theorem}
证明将在紧接的两小节完成。这里最后一句来自约数一总在求和范围内，故表示数严格为正；不能把所有项为非负误当成已知总和不为零。该定理称为\term[jacobi-sipingfanghe]{Jacobi 四平方和定理}[Jacobi four-square theorem]。

\subsection{\texorpdfstring{\(\theta\)}{\AnalysisPdfTheta} 函数证明}
\label{b3:subsec:240403}

取格 \(\Lambda=\pi\mathbb Z+\pi\tau\mathbb Z\)，记 \(e_1=\wp(\pi/2)\)、\(e_3=\wp(\pi\tau/2)\)。由准周期性，\(-\mathrm{d}^2\log\theta_1(z)/\mathrm{d}z^2\) 双周期；在格点主部为 \(z^{-2}\)。与 \(\wp\) 相减是整椭圆函数，所以
\(\wp(z)=-(\log\theta_1(z))''+C\)。平移级数得 \(\theta_1(z+\pi/2)=\theta_2(z)\)，以及
\(\theta_1(z+\pi\tau/2)=\mathrm{i}q^{-1/4}e^{-\mathrm{i}z}\theta_4(z)\)。三个偶函数在原点导数为零，因此
\[
e_1-e_3=\frac{\theta_4''(0)}{\theta_4}-\frac{\theta_2''(0)}{\theta_2}.
\]
另一方面，函数
\(F(z)=(\theta_1'(0)/\theta_2)^2(\theta_2(z)/\theta_1(z))^2\) 双周期，在原点主部为 \(z^{-2}\)，并在 \(\pi/2\) 有二重零点；故 \(F=\wp-e_1\)。再代入 \(z=\pi\tau/2\)，利用
\(\theta_2(z+\pi\tau/2)=q^{-1/4}e^{-\mathrm{i}z}\theta_3(z)\) 和\eqref{b3:eq:24-derivative-identity}，得到 \(e_3-e_1=-\theta_3^4\)。于是
\[
\frac{\theta_4''(0)}{\theta_4}-\frac{\theta_2''(0)}{\theta_2}=\theta_3^4.
\]
每个 theta 级数还满足 \(q\partial_q\theta_j=-\tfrac14\partial_z^2\theta_j\)，直接逐项微分即可核对。因此
\begin{equation}\label{b3:eq:24-log-derivative}
4q\frac{\mathrm{d}}{\mathrm{d}q}\log(\theta_2/\theta_4)=\theta_3^4.
\end{equation}
整个推导中的除法由三重积保证分母非零，椭圆函数的比较则只比较同一格下的主部和一个零点，故不存在未确定的比例常数。

\subsection{生成函数方法}
\label{b3:subsec:240404}

\begin{proof}[\cref{b3:thm:24-four-squares}的证明]
对 \(\theta_2/\theta_4\) 的乘积取对数导数，由\eqref{b3:eq:24-log-derivative}得
\[
\theta_3^4=1+16\sum_{n\ge1}\frac{nq^{2n}}{1+q^{2n}}
+8\sum_{m\ge1,\ m\text{ odd}}\frac{mq^m}{1-q^m}.
\]
恒等式 \(x/(1+x)=x/(1-x)-2x^2/(1-x^2)\) 将偶数部分改写为
\(8\sum_{d\text{ even}}d q^d/(1-q^d)-32\sum_{n\ge1}nq^{4n}/(1-q^{4n})\)。与奇数部分相合，即为
\[
\theta_3^4=1+8\sum_{d\ge1,\ 4\nmid d}\frac{dq^d}{1-q^d}.
\]
在 \(|q|\le r<1\) 上所有对数导数级数和几何展开都绝对一致收敛，因为 \(\sum d r^d/(1-r^d)<\infty\)。故可把 \(q^d/(1-q^d)=\sum_{m\ge1}q^{dm}\) 代入并收集系数，得到所述约数和。最后与\secref{b3:sec:2403}中的生成函数比较。
\end{proof}
例如奇素数 \(p\) 的四平方和表示数为 \(8(p+1)\)；\(n=4\) 时为 \(8(1+2)=24\)，包括八个仅一个坐标为 \(\pm2\) 的表示与十六个四坐标均为 \(\pm1\) 的表示。这个核算检验了符号、顺序及排除四倍约数的三个约定。

\ProseParagraphBreak

半周期比较中的主部与常数值得再逐点核验。由于 \(\theta_1\) 在0单零，\(\theta_1(z)=\theta_1'(0)z+O(z^3)\)；偶函数 \(\theta_2(z)=\theta_2+O(z^2)\)。因此
\((\theta_1'(0)/\theta_2)^2(\theta_2(z)/\theta_1(z))^2=z^{-2}+O(1)\)，没有一阶极项。平方后的商在两周期下的符号和指数相消，确为椭圆函数。它与 \(\wp\) 的差因主部相同而整，故为常数；在 \(z=\pi/2\) 时分子为 \(\theta_1(\pi)=0\)，分母为 \(\theta_2\ne0\)，常数恰为 \(-e_1\)。这一步同时用到了零点的位置、单重性及常数非零，三者均由前面的级数或乘积给出。
\ProseParagraphBreak
在另一半周期，平移公式逐项给
\(\theta_2(\pi\tau/2)=q^{-1/4}\theta_3\) 与
\(\theta_1(\pi\tau/2)=\mathrm{i}q^{-1/4}\theta_4\)。代入平方商，先得到比值平方 \(-\theta_3^2/\theta_4^2\)，再乘
\((\theta_1'(0)/\theta_2)^2=\theta_3^2\theta_4^2\)，最后才得到 \(-\theta_3^4\)。负号来自 \(\mathrm{i}^2=-1\)，不是法向或周期方向的约定；四次方来自两个平方因子的相乘。明确这些步骤后，热方程的负四分之一因子也与最后的正号一致。
\ProseParagraphBreak
从乘积到 Lambert 级数同样可逐类追踪。\(\theta_2/\theta_4\) 中公共的 \(1-q^{2n}\) 全部消去，剩下常数 \(2q^{1/4}\)、偶幂的 \((1+q^{2n})^2\) 和奇幂的 \((1-q^{2n-1})^{-2}\)。乘上 \(4q\mathrm{d}/\mathrm{d}q\) 后，第一项给一，偶幂项给 \(16nq^{2n}/(1+q^{2n})\)，奇幂项给 \(8(2n-1)q^{2n-1}/(1-q^{2n-1})\)。把偶数项改写后，其中全部四的倍数恰被减掉一次。若忘记平方因子或 \(q^{1/4}\)，会分别造成表示数少一倍或漏掉 \(r_4(0)=1\)；前几项计数可以迅速检查这些错误。
\ProseParagraphBreak
例如 \(r_4(2)=24\)：选择两个非零位置有六种，各自符号有四种。\(r_4(3)=32\)：选择零坐标有四种，其余三个符号有八种。对 \(n=9\)，一类表示为 \((\pm3,0,0,0)\) 的排列，共八个；另一类是 \((\pm2,\pm2,\pm1,0)\)，零位置四种、绝对值一的位置三种、三个非零符号八种，共九十六个。总数一百零四，正好是 \(8(1+3+9)\)。相同绝对值坐标的置换不能再重复乘一个阶乘，否则会过计。
\ProseParagraphBreak
四平方公式还能给比体积法更细的累计误差。设 \(X\ge1\)，交换有限约数和得到
\[
 \sum_{1\le n\le X}r_4(n)=8\sum_{m\le X}\sum_{d\le X/m,\,4\nmid d}d.
\]
由 \(\sum_{d\le Y}d=Y^2/2+O(Y)\)，扣去四的倍数后内和为 \(3Y^2/8+O(Y)\)。于是总和为
\[
 3X^2\sum_{m\le X}m^{-2}+O\left(X\sum_{m\le X}m^{-1}\right)
 =\frac{\pi^2}{2}X^2+O(X\log(X+2)).
\]
最后使用正弦乘积比较所得 \(\sum m^{-2}=\pi^2/6\)，尾和为 \(O(1/X)\)。主项正是四维单位球体积 \(\pi^2/2\)，因此复分析推导的精确系数与实几何计数相符。此处取的是累计和；不能把其误差逐项作差就得到每个 \(r_4(n)\) 的同等精度。

% !TeX root = ../../Book3.tex
\OptionalSection{二次型的 \texorpdfstring{\(\theta\)}{\AnalysisPdfTheta} 级数}{b3:sec:2405}

把平方和换成一般正定二次型，收敛和生成函数意义仍保留，但参数反演会涉及对偶格，模群也通常缩小。

\subsection{正定二次型}
\label{b3:subsec:240501}

设 \(A\) 为实对称正定矩阵，\(Q(x)=\tfrac12x^TAx\)。正定性给 \(Q(x)\ge c|x|^2\)，因而 \(\sum_{m\in\mathbb Z^d}e^{2\pi\mathrm{i}\tau Q(m)}\) 在上半平面正规收敛。若 \(A\) 为整矩阵且对角线为偶数，则 \(Q(m)\) 为整数，级数对 \(\tau\mapsto\tau+1\) 不变。

\subsection{格 \texorpdfstring{\(\theta\)}{\AnalysisPdfTheta} 级数}
\label{b3:subsec:240502}

对欧氏格 \(L\)，记 \(L^*=\{y:\langle x,y\rangle\in\mathbb Z\ (x\in L)\}\)，基本区域体积为 \(v(L)\)。格 theta 级数为
\(\Theta_L(\tau)=\sum_{x\in L}e^{\pi\mathrm{i}\tau|x|^2}\)。把实高斯的 Fourier 变换按基矩阵换元，Poisson 求和给
\[
\Theta_L(-1/\tau)=\frac{(-\mathrm{i}\tau)^{d/2}}{v(L)}\Theta_{L^*}(\tau).
\]
高维求和公式可先对标准整数格逐个坐标作周期 Fourier 展开，再线性换元；高斯的共同快速衰减保证所有换序。

\subsection{模性}
\label{b3:subsec:240503}

若 \(L=L^*\)，并且每个格向量平方长度为偶整数，则 \(v(L)=1\)，平移与反演都封闭在同一个 theta 级数上。若另有 \(d\) 是八的倍数，则 \((-\mathrm{i})^{d/2}=1\)，两个生成元给出权 \(d/2\) 的全模群变换律；尖点常数项一及指数尾部给尖点全纯性。因此 \(\Theta_L\in M_{d/2}\)。这里把维数条件列为假设，并不在本节证明偶自对偶格的维数必为八的倍数。

\subsection{表示数}
\label{b3:subsec:240504}

若 \(Q\) 在整格上取非负整数值，则 Fourier 系数等于表示数 \(r_Q(n)=\#\{m:Q(m)=n\}\)。正定性保证每层有限。对非自对偶格，反演把它送到对偶格或若干陪集级数的组合，通常得到子群上的模形式或向量值变换律；不能从一个高斯公式直接宣称所有整二次型都给全模群标量模形式。\appref{b3:app:E}将说明有限层级的相应语言。

\ProseParagraphBreak

高维 Gaussian 公式中的体积因子可在矩阵形式下直接看到。若 \(A\) 正定，选正交坐标将它对角化，再逐坐标使用一维 Gaussian 变换，得到
\[
 \widehat{e^{-\pi t x^TAx}}(\xi)
 =t^{-d/2}(\det A)^{-1/2}e^{-\pi\xi^TA^{-1}\xi/t}.
\]
正定性确保特征值全正、Gaussian 可积；逆矩阵出现在频率的二次型中，行列式来自实变量换元。因此反演自然交换原格与对偶格，而不会对任意正定整矩阵保持同一个级数。若二次型只是半正定，例如 \(Q(x,y)=x^2\)，每个固定表示层中 \(y\) 可任意变化，系数已经无穷；若不定，例如 \(x^2-y^2\)，零层也无穷。正定性不仅用于估计，还保证生成函数本身的计数意义。
\ProseParagraphBreak
偶自对偶条件有一个具体的八维实例。令
\[
 D_8=\{x\in\mathbb Z^8:\textstyle\sum_jx_j\text{ 为偶数}\},\qquad
 L=D_8\cup\left(D_8+\left(\tfrac12,\ldots,\tfrac12\right)\right).
\]
两次半整数平移回到 \(D_8\)，故并集是加法群。整数部分的平方长度模二等于坐标和，因而偶；半整数部分 \(x+v\) 的平方长度为 \(|x|^2+\sum x_j+2\)，也为偶整数。由极化恒等式，任意两格向量的内积为整数，所以 \(L\subset L^*\)。\(D_8\) 在整数格中指标为二，体积二；\(L\) 又含它为指标二子格，故体积一。对偶格体积为倒数一，包含关系与相同体积迫使 \(L=L^*\)。这给出本节模性结论所需的全部假设。
\ProseParagraphBreak
还可独立数出其最短非零向量。整数部分长度平方为二的向量恰有两个坐标为 \(\pm1\)，共 \(\binom82\cdot4=112\) 个。半整数部分若有坐标绝对值至少 \(3/2\)，平方长度已至少四；所以长度平方二时八个坐标全为 \(\pm1/2\)。属于指定陪集要求负号数为偶数，恰有 \(2^7=128\) 个。总数为240，与权四 Eisenstein 级数的一次系数相符。这个几何核算检验了长度平方 \(2n\) 与模变量 \(e^{2\pi\mathrm{i}\tau}\) 中 \(q^n\) 的对应，也给抽象的偶自对偶条件一个可以逐坐标检查的模型。

% 本章小结（不计入编号）
\begin{chaptersummary}
二平方和的障碍来自素因子在 Gaussian 整数中的不同分解方式；四平方和的精确数量则来自 theta 乘积的对数导数。后者的关键连接是：先用乘积确定单零点和 Jacobi 导数常数，再比较同一空间格上的椭圆函数，在两个半周期取值，最后用热方程转成参数导数。这样得到的 Lambert 级数可以按约数整理，正文证明不依赖章末练习。

级数、乘积、空间格和模形式各自携带归一化。尤其本章 \(q=e^{\pi\mathrm{i}\tau}\)，而偶格的自然计数变量为 \(q_{\mathrm{mod}}=q^2\)。一般格始终有包含对偶格和基本区域体积的反演公式；只有进一步核对偶性、自对偶性及权数后，才能直接使用前章全模群模形式的维数结论。累计格点数的体积近似又说明，同一表示数既有解析结构，也有几何意义，但累计估计不能替代单个系数的精确公式。
\end{chaptersummary}

\BookExercises

% 八题重新组织：前六题使用基本路线，后两题使用选读格理论。
% 二平方和、三角数、热核、反演固定点、椭圆退化和倍参数恒等式均由正文工具推导；
% 末两题考察模变换的条件与模形式有限维性的应用，不承担正文证明依赖。
\begin{exercise}\label{b3:ex:24-1}
对正整数 \(n\)，令
\[
P_2(n)=\#\{(x,y)\in\mathbb Z^2:x^2+y^2=n,\ 
                                  \gcd(|x|,|y|)=1\}.
\]
规定 \(\gcd(a,0)=a\)（\(a>0\)），不同次序和符号分别计数。求 \(P_2(n)\) 的素因子公式，并比较 \(P_2(25)\) 与 \(r_2(25)\)。
\end{exercise}
\begin{solution}
若 \(n\) 含模四为三的素因子 \(\ell\)，则其指数为奇数时没有任何表示；指数为正偶数 \(2s\) 时，正文的分解公式使每个 \(\alpha=x+\mathrm{i}y\) 都含实因子 \(\ell^s\)，从而 \(\ell\mid x,y\)。所以原始表示中完全不能出现这种素因子。

写剩余的分解为
\[
n=2^a\prod_{j=1}^k p_j^{a_j},
\qquad p_j\equiv1\pmod4,\quad a_j\ge1.
\]
正文已证 \(\alpha\) 中 \(1+\mathrm{i}\) 的指数恰为 \(a\)。若 \(a\ge2\)，由于 \((1+\mathrm{i})^2=2\mathrm{i}\)，两个坐标均为偶数；若 \(a\le1\)，则二不能同时整除两个坐标。

对每个 \(p_j=\pi_j\bar\pi_j\)，\(\alpha\) 中的因子为
\(\pi_j^{r_j}\bar\pi_j^{a_j-r_j}\)。
实整数 \(p_j\) 整除 \(\alpha\)，当且仅当两个指数都至少为一，这又等价于 \(p_j\) 同时整除 \(x,y\)。故原始性要求
\(r_j=0\) 或 \(r_j=a_j\)，每个不同奇素因子给两种选择。反过来，若这些条件都满足，则任何可能同时整除两坐标的有理素数都已被排除，故最大公因子为一。四个单位分别计数，得到
\[
P_2(n)=
\begin{cases}
4\cdot2^k,&n=2^a\prod_{j=1}^k p_j^{a_j},\
                 a\in\{0,1\},\ p_j\equiv1\pmod4,\\
0,&\text{其余情形}.
\end{cases}
\]
空乘积允许 \(n=1,2\)，此时均有四个原始表示。对 \(25\)，
\(P_2(25)=8\)，对应绝对值型 \((3,4)\)；而 \(r_2(25)=12\)，另外四个为 \((\pm5,0),(0,\pm5)\)，它们的坐标有公因子五。
\end{solution}

\begin{exercise}\label{b3:ex:24-2}
令 \(T(a)=a(a+1)/2\)（\(a\in\mathbb Z_{\ge0}\)）。对整数 \(n\ge0\)，用 \(t_4(n)\) 表示满足
\(T(a_1)+T(a_2)+T(a_3)+T(a_4)=n\) 的有序非负整数四元组数。证明
\[
t_4(n)=\sigma_1(2n+1),
\qquad \sigma_1(m)=\sum_{\substack{d\mid m\\d>0}}d.
\]
说明这如何推出每个非负整数都是四个非负三角数之和。
\end{exercise}
\begin{solution}
恒等式 \(8T(a)+1=(2a+1)^2\) 将原问题化为
\[
(2a_1+1)^2+\cdots+(2a_4+1)^2=8n+4.
\]
令 \(m=2n+1\)，则 \(m\) 为正奇数。四个平方之和为 \(4m\) 时，平方模四的值为零或一，所以奇坐标的个数只能为零或四。因此全部表示恰分为全偶和全奇两类。全偶表示除以二，与 \(m\) 的四平方和表示一一对应，数量为 \(r_4(m)\)；全奇表示数就是
\(r_4(4m)-r_4(m)\)。

因为 \(m\) 为奇数，Jacobi 公式给
\[
r_4(m)=8\sigma_1(m),\qquad
r_4(4m)=8\!\sum_{d\mid m}(d+2d)=24\sigma_1(m).
\]
后式中允许的约数只有 \(d,2d\)，所有 \(4d\) 被排除，故全奇表示共有 \(16\sigma_1(m)\) 个。

每个全奇表示的四个坐标都非零。独立改变四个符号给恰好十六个不同表示，而取各坐标绝对值以后，
\(a_j=(|x_j|-1)/2\) 唯一确定。除以十六，便得所求公式。由于 \(\sigma_1(2n+1)\ge1\)，表示总存在；\(n=0\) 时公式给一，正对应四个 \(a_j\) 都为零。
\end{solution}

\begin{exercise}\label{b3:ex:24-3}
对 \(t>0\)、\(x\in\mathbb R\)，令
\[
K_t(x)=\frac1\pi\theta_3(x,\mathrm{i}t).
\]
证明 \(K_t\) 为正的 \(\pi\)-周期函数，
\(\int_0^\pi K_t(x)\,\mathrm dx=1\)，并满足
\[
\int_0^\pi K_s(x-y)K_t(y)\,\mathrm dy=K_{s+t}(x)
\quad(s,t>0).
\]
若 \(f\) 为连续的 \(\pi\)-周期函数，证明
\(\int_0^\pi K_t(x-y)f(y)\,\mathrm dy\) 在 \(t\downarrow0\) 时一致趋于 \(f(x)\)。
\end{exercise}
\begin{solution}
在高斯 Poisson 公式中把实参数换成 \(1/t\)，取平移参数 \(a=-x/\pi\)、调制参数 \(b=0\)，再用 \(\theta_3\) 关于空间变量的偶性，得
\[
K_t(x)=\frac1{\pi\sqrt t}
       \sum_{m\in\mathbb Z}
       \exp\!\left(-\frac{(x-\pi m)^2}{\pi t}\right).
\]
每一项都正，平移 \(x\mapsto x+\pi\) 只重排指标。对非负级数逐项积分，将各个长度为 \(\pi\) 的区间拼成实轴，得到
\[
\int_0^\pi K_t(x)\,\mathrm dx
=\frac1{\pi\sqrt t}\int_{\mathbb R}e^{-u^2/(\pi t)}\,\mathrm du=1.
\]
这也解释了前因子 \(1/\pi\)：原来的 theta 函数在一个空间周期上的积分为 \(\pi\)。

由原级数
\[
K_t(x)=\frac1\pi\sum_{n\in\mathbb Z}
               e^{-\pi n^2t}e^{2\mathrm{i}nx}.
\]
当 \(s,t>0\) 固定时，两个级数绝对一致收敛，其乘积可逐项积分。利用
\(\int_0^\pi e^{2\mathrm{i}(m-n)y}\,\mathrm dy
=\pi\)（\(m=n\)）或零（\(m\ne n\)），只留下相同频率，系数相乘成为
\(e^{-\pi n^2(s+t)}\)，即所述卷积公式。

最后把一个周期取为 \([-\pi/2,\pi/2]\)。对 \(0<\delta<\pi/2\)，高斯表达式给
\[
\int_{\delta\le |y|\le\pi/2}K_t(y)\,\mathrm dy
\le \frac1{\pi\sqrt t}
             \int_{|u|\ge\delta}e^{-u^2/(\pi t)}\,\mathrm du
\le \frac{\sqrt t}{\delta}e^{-\delta^2/(\pi t)}.
\]
第一个不等式把各平移区间合并，并扩大到 \(|u|\ge\delta\)；
第二个用 \(u/\delta\ge1\) 估计正半轴积分，再利用对称性。

连续周期函数在实轴上一致连续。给定 \(\varepsilon>0\)，可先选 \(\delta\)，使
\(|f(x-y)-f(x)|<\varepsilon\) 对所有 \(x\) 及 \(|y|<\delta\) 成立。利用质量为一，卷积与 \(f(x)\) 之差的绝对值至多为
\[
\varepsilon+
2\|f\|_\infty\frac{\sqrt t}{\delta}e^{-\delta^2/(\pi t)}.
\]
右侧与 \(x\) 无关，再令 \(t\downarrow0\)，即得一致收敛。这同时给出了本章热方程中周期初值的一个具体构造。
\end{solution}

\begin{exercise}\label{b3:ex:24-4}
设 \(T(t)=\sum_{n\in\mathbb Z}e^{-\pi n^2t}\)（\(t>0\)）。
利用 \(T(t)=t^{-1/2}T(1/t)\) 证明
\[
\sum_{n\in\mathbb Z}n^2e^{-\pi n^2}
=\frac1{4\pi}\sum_{n\in\mathbb Z}e^{-\pi n^2}.
\]
再说明函数 \(H(s)=e^{s/4}T(e^s)\)（\(s\in\mathbb R\)）为何是偶函数，并用它解释上式中的固定参数 \(t=1\)。
\end{exercise}
\begin{solution}
在包含一的紧区间上，导数级数由常数倍的
\(\sum n^2e^{-c n^2}\) 控制，故可逐项求导。
反演式微分为
\[
T'(t)=-\tfrac12t^{-3/2}T(1/t)-t^{-5/2}T'(1/t).
\]
取 \(t=1\)，得到 \(2T'(1)=-T(1)/2\)。
代入 \(T'(1)=-\pi\sum n^2e^{-\pi n^2}\)，就是所求恒等式。

另一方面，反演式在 \(t=e^s\) 处给
\(T(e^s)=e^{-s/2}T(e^{-s})\)，因此
\[
H(s)=e^{-s/4}T(e^{-s})=H(-s).
\]
参数反演 \(t\mapsto1/t\) 在对数坐标中成为反射
\(s\mapsto-s\)，其固定点为 \(s=0\)，即 \(t=1\)。
偶性要求 \(H'(0)=0\)，而
\(H'(0)=T(1)/4+T'(1)\)，再次给出同一矩恒等式。
\end{solution}

\begin{exercise}\label{b3:ex:24-5}
设 \(\tau\in\mathfrak H\)、\(q=e^{\pi\mathrm{i}\tau}\)，
\(\Lambda_\tau=\pi\mathbb Z+\pi\tau\mathbb Z\)。
在 \(|\operatorname{Im}z|<\pi\operatorname{Im}\tau\)、\(z\notin\pi\mathbb Z\) 内证明
\[
\wp(z;\Lambda_\tau)=\frac1{\sin^2z}-\frac13
 +8\sum_{m\ge1}\frac{m q^{2m}}{1-q^{2m}}
                         \bigl(1-\cos(2mz)\bigr).
\]
进而证明，当 \(t\to\infty\) 时，
\[
\wp(z;\pi\mathbb Z+\pi\mathrm{i}t\mathbb Z)
\longrightarrow\frac1{\sin^2z}-\frac13
\]
在 \(\mathbb C\setminus\pi\mathbb Z\) 的每个紧集上一致成立。
\end{exercise}
\begin{solution}
正文的对数导数公式在所述条带内为
\[
\frac{\theta_1'(z)}{\theta_1(z)}
=\cot z+4\sum_{m\ge1}
              \frac{q^{2m}}{1-q^{2m}}\sin(2mz).
\]
在条带的紧子集上，各阶导数都受带多项式因子的几何级数控制，所以可以求导，也可以在原点比较 Laurent 展开的系数。
由 \(\cot z=z^{-1}-z/3+O(z^3)\)，上式中 \(z\) 的系数为
\[
-\frac13+8\sum_{m\ge1}\frac{m q^{2m}}{1-q^{2m}}.
\]
另一方面，正文展开
\(\theta_1(z)=Az+Bz^3+O(z^5)\) 给出该系数为
\(2B/A=\theta_1'''(0)/(3\theta_1'(0))\)。
将此值代入\cref{b3:prop:24-theta-wp}，并对上面的对数导数求导，便得到所求公式。常数 \(-1/3\) 正是使原点 Laurent 展开的常数项为零所必需的修正。

设紧集 \(K\subset\mathbb C\setminus\pi\mathbb Z\)，并令
\(R=\max_{z\in K}|\operatorname{Im}z|\)。当 \(t>R/\pi\) 时，\(K\) 位于公式的条带内。取 \(\tau=\mathrm{i}t\)，则
\(q=e^{-\pi t}\)，且
\[
|1-\cos(2mz)|\le1+e^{2mR}\quad(z\in K).
\]
对于充分大的 \(t\)，分母 \(1-e^{-2\pi mt}\) 一致远离零，所以级数尾项的总和由
\[
C\sum_{m\ge1}m e^{-2\pi mt}(1+e^{2mR})
=O_K(e^{-2\pi t})
\]
控制。这给出紧集上一致收敛，甚至给出指数误差。几何上，第二个周期向虚方向变长，有限区域内留下的极点排列趋于实轴上的 \(\pi\mathbb Z\)。
\end{solution}

\begin{exercise}\label{b3:ex:24-6}
对 \(\tau\in\mathfrak H\)，记 \(\theta_j(\tau)=\theta_j(0,\tau)\)。
直接重排绝对收敛的双重级数，证明倍参数恒等式
\[
\begin{aligned}
\theta_3(\tau)^2&=\theta_3(2\tau)^2+\theta_2(2\tau)^2,\\
\theta_4(\tau)^2&=\theta_3(2\tau)^2-\theta_2(2\tau)^2,\\
\theta_2(\tau)^2&=2\theta_2(2\tau)\theta_3(2\tau).
\end{aligned}
\]
据此另行推出 \(\theta_3(\tau)^4=\theta_2(\tau)^4+\theta_4(\tau)^4\)。
\end{exercise}
\begin{solution}
在 \(\theta_3(\tau)^2\) 中，令
\(u=m+n,\ v=m-n\)。整数对 \((m,n)\) 与同奇偶的整数对 \((u,v)\) 一一对应，逆变换为
\(m=(u+v)/2,\ n=(u-v)/2\)，且
\(m^2+n^2=(u^2+v^2)/2\)。
当 \(u=2a,v=2b\) 时，贡献为
\(\sum_{a,b}q^{2a^2+2b^2}=\theta_3(2\tau)^2\)；
当 \(u=2a+1,v=2b+1\) 时，贡献为
\(\sum_{a,b}q^{2(a+1/2)^2+2(b+1/2)^2}=\theta_2(2\tau)^2\)。
这证明第一式。

\(\theta_4(\tau)^2\) 多出符号
\((-1)^{m+n}=(-1)^u\)，在全偶部分为正、全奇部分为负，故得到第二式。
对 \(\theta_2(\tau)^2\)，改用
\[
u=(m+\tfrac12)+(n+\tfrac12)=m+n+1,\qquad
v=(m+\tfrac12)-(n+\tfrac12)=m-n.
\]
此时 \(u,v\) 为异奇偶整数，逆变换仍唯一。
偶奇和奇偶两部分各给一个
\(\theta_3(2\tau)\theta_2(2\tau)\)，于是得到第三式。
所有这些分组都由高斯型双重级数的绝对收敛保证，不存在条件收敛下的重排问题。

最后将前两式相减、相加，再相乘：
\[
\theta_3(\tau)^4-\theta_4(\tau)^4
=4\theta_2(2\tau)^2\theta_3(2\tau)^2
=\theta_2(\tau)^4.
\]
这提供了四次恒等式的另一条路线，并说明倍参数公式与正文的奇偶格点变换有共同的计数基础。
\end{solution}

\begin{optionalexercise}\label{b3:ex:24-7}
在标准欧氏平面中取
\[
L=\operatorname{diag}(\sqrt2,1/\sqrt2)\mathbb Z^2.
\]
求其基本区域体积与对偶格，证明 \(L\ne L^*\)，但
\(\Theta_L=\Theta_{L^*}\)。写出 \(\Theta_L(-1/\tau)\) 的变换律，并证明 \(\Theta_L\) 不满足 \(\tau\mapsto\tau+1\) 的不变性。由此解释，为何体积为一，甚至反演公式在同一个函数上闭合，仍不足以得到全模群模形式。
\end{optionalexercise}
\begin{solution}
基矩阵的行列式为一，所以 \(v(L)=1\)，而逆转置给
\[
L^*=\operatorname{diag}(1/\sqrt2,\sqrt2)\mathbb Z^2.
\]
向量 \((0,1/\sqrt2)\) 属于 \(L\) 而不属于 \(L^*\)，故两格不同；交换两个坐标却将 \(L\) 等距映到 \(L^*\)，所以长度计数相同。由定义，
\[
\Theta_L(\tau)
=\sum_{m,n\in\mathbb Z}e^{\pi\mathrm{i}\tau(2m^2+n^2/2)}
=\theta_3(0,2\tau)\theta_3(0,\tau/2).
\]
一般格反演在维数二给
\(\Theta_L(-1/\tau)=(-\mathrm{i}\tau)\Theta_{L^*}(\tau)
=-\mathrm{i}\tau\Theta_L(\tau)\)。

但 \(L\) 的最短非零向量为 \((0,\pm1/\sqrt2)\)，平方长度是 \(1/2\)，所以当 \(\operatorname{Im}\tau\to\infty\) 时，
\[
\Theta_L(\tau)=1+2e^{\pi\mathrm{i}\tau/2}
                 +O(e^{-2\pi\operatorname{Im}\tau}).
\]
这个估计对 \(\operatorname{Re}\tau=0,1\) 均成立，由绝对收敛级数中剩余平方长度至少为二得到。
平移 \(\tau\mapsto\tau+1\) 将显示出的第二项乘以 \(\mathrm{i}\)，因而
\[
\Theta_L(1+\mathrm{i}t)-\Theta_L(\mathrm{i}t)
=2(\mathrm{i}-1)e^{-\pi t/2}+O(e^{-2\pi t})\ne0
\]
对充分大的 \(t\) 成立。
这里体积为一没有保证自对偶，也没有保证长度平方为偶整数。反演中的额外常数 \(-\mathrm{i}\) 和平移失败都表明，不能直接套用前章全模群的整数权变换律。
\end{solution}

\begin{optionalexercise}\label{b3:ex:24-8}
设 \(L\) 是标准欧氏空间中的正定偶自对偶格，记
\[
R_L(n)=\#\{v\in L:|v|^2=2n\},\qquad
Q=e^{2\pi\mathrm{i}\tau}.
\]
当 \(\dim L=16\) 时，求 \(R_L(1)\) 与 \(R_L(2)\)。
当 \(\dim L=24\) 时，用 \(R_L(1)\) 表示整个 theta 级数，并证明
\(R_L(2)=196560-24R_L(1)\)。
特别地，若某个二十四维偶自对偶格没有平方长度为二的向量，它有多少个平方长度为四的向量？本题不要求证明这样的格存在。
\end{optionalexercise}
\begin{solution}
末节已证 \(\Theta_L\in M_{\dim L/2}\)，而正定性使常数项为一，偶性使展开为
\(\Theta_L=1+\sum_{n\ge1}R_L(n)Q^n\)。
这里使用 \(Q=q^2\)，所以 \(Q^n\) 计平方长度 \(2n\)，不能读成平方长度 \(n\)。

由\cref{b3:thm:23-modular-ring}，\(M_8=\mathbb C E_4^2\)。
十六维情形比较常数项便有 \(\Theta_L=E_4^2\)。
又由\eqref{b3:eq:23-e4-e6-fourier}，
\[
E_4=1+240Q+2160Q^2+O(Q^3).
\]
平方后得
\[
E_4^2=1+480Q+61920Q^2+O(Q^3).
\]
因此 \(R_L(1)=480,\ R_L(2)=61920\)，这些数量与具体选择哪个十六维偶自对偶格无关。

二十四维时，\(M_{12}=\mathbb C E_4^3\oplus\mathbb C\Delta\)。
常数项一先确定 \(E_4^3\) 的系数为一，再比较一次项，得到
\[
\Theta_L=E_4^3+(R_L(1)-720)\Delta.
\]
为求第二个系数，使用正文已算出的
\[
E_4^3=1+720Q+179280Q^2+O(Q^3),\qquad
\Delta=Q-24Q^2+O(Q^3).
\]
于是
\[
R_L(2)=179280-24(R_L(1)-720)
       =196560-24R_L(1).
\]
若 \(R_L(1)=0\)，就有恰好 \(196560\) 个平方长度为四的向量。这个推论从给定格的条件出发；它计算可能的长度分布，并不构造或分类这样的格。
\end{solution}
